\section{Summary of Claims and Status}
%=============================================================================

\begingroup
\small
\renewcommand{\arraystretch}{1.15}
\begin{longtable}{@{}p{3cm}p{4cm}p{3cm}p{5.2cm}@{}}
\caption{TFPT status dashboard. ``Exact'': $<0.1\%$; ``Solid'': $<1\%$; ``Minor'': $<3\%$; ``Moderate'': $<5\%$.}\label{tab:status}\\
\toprule
\textbf{Domain} & \textbf{Quantity} & \textbf{Deviation} & \textbf{Status}\\
\midrule
\endfirsthead
\toprule
\textbf{Domain} & \textbf{Quantity} & \textbf{Deviation} & \textbf{Status}\\
\midrule
\endhead
\midrule
\multicolumn{4}{r}{\small Continued on next page.}\\
\midrule
\endfoot
\bottomrule
\endlastfoot
\multicolumn{4}{l}{\textit{Structural (parameter-free)}}\\
\quad CFE & $\alpha^{-1}(0)=\num[round-mode=places,round-precision=12]{\TFPTalphaInvGOverFourNum}$ & $z=\num[round-mode=places,round-precision=3]{\TFPTalphaZGOverFourNum}$ & \textcolor{green!60!black}{\textbf{Derived (double cover + $g=5$ defect partition; no fits)}}\\
\quad UFE & $\beta = \num[round-mode=places,round-precision=4]{\TFPTbetaDegNum}^\circ$ & 0.5--1.7$\sigma$ & \textcolor{green!60!black}{\textbf{Solid}}\\
\quad Cabibbo & $\lambda = \num[round-mode=places,round-precision=6]{0.2243762369}$ & $+0.03\%$ & \textcolor{green!60!black}{\textbf{Solid}}\\
\midrule
\multicolumn{4}{l}{\textit{RG-dressed}}\\
\quad Two-loop & $\alpha_3(1\,\text{PeV})$ & $-1.38\%$ & \textcolor{blue}{\textbf{Good}}\\
\quad Two-loop & $\alpha_3(\mu_{\cthree})$ & $0.00\%$ & \textcolor{green!60!black}{\textbf{Exact}}\\
\quad Unification gate & $\alpha_1,\alpha_2,\alpha_3$ & mismatch $\sim 0.38\%$ & \textcolor{green!60!black}{\textbf{Solid (explicit gate; discrete policy scan)}}\\
\midrule
\multicolumn{4}{l}{\textit{Structural Extensions (Part II)}}\\
\quad PMNS ($\theta_{13}$) & $\sin^2\theta_{13} = \num[round-mode=places,round-precision=5]{0.0231084352}$ & $+5.04\%$ & \textcolor{blue}{\textbf{Testable}}\\
\quad PMNS ($\theta_{12}$) & $\sin^2\theta_{12} = \num[round-mode=places,round-precision=3]{0.318}$ & $+3.56\%$ & \textcolor{blue}{\textbf{Derived}}\\
\quad Inflation ($n_s$) & $n_s \approx 0.9643$ ($N\approx 56$) & within errors & \textcolor{blue}{\textbf{Derived (K1--K3)}}\\
\quad Inflation ($r$) & $r \approx 0.0038$ & future & \textcolor{blue}{\textbf{Derived (K1--K3)}}\\
\quad Bounce transfer function & $T(k)$ large-scale features & future & \textcolor{orange}{\textbf{Conditional (bounce module)}}\\
\quad Axion DM & $m_a \approx 64\,\mu\text{eV}$ & $\theta_i$ model & \textcolor{orange}{\textbf{Structured}}\\
\midrule
\multicolumn{4}{l}{\textit{Suite-closure (physics-mode; assumption-explicit proxy modules)}}\\
\quad Global consistency (suite, physics) & $\chi^2$ scorecard & $p\simeq 0.435$ ($\chi^2\simeq 15.23$, 15 dof) & \textcolor{green!60!black}{\textbf{Consistent (dashboard; no full likelihood)}}\\
\quad Flavor (suite, physics) & CKM/PMNS $\chi^2$ & $p_{\mathrm{CKM}}\simeq 0.242$, $p_{\mathrm{PMNS}}\simeq 0.525$ & \textcolor{green!60!black}{\textbf{Consistent (imported $\chi^2$ terms)}}\\
\quad Axion DM (relic) & $\Omega_a h^2 \approx 0.123$ & $+2.2\%$ & \textcolor{blue}{\textbf{Derived (post-inflation; $C_{\mathrm{str}}=7/3$)}}\\
\quad Dark energy (norm) & $\rho_\Lambda$ via $e^{-\alpha^{-1}/2}$ & 0.086 dex & \textcolor{green!60!black}{\textbf{Candidate (theoryv3 audit; $n=1/2$)}}\\
\quad Dark energy (condensate) & $\rho_\Lambda$ torsion gap equation & $z\simeq 2.14$ (0.215 dex) & \textcolor{orange}{\textbf{Close, but not closed (open interface)}}\\
\quad BBN & $Y_p, D/H, N_\mathrm{eff}$ & within $2\sigma$ & \textcolor{blue}{\textbf{Consistent (suite)}}\\
\quad Baryogenesis & $\eta_b \approx 5.8\times10^{-10}$ & within 0.5 dex & \textcolor{blue}{\textbf{Proxy (suite)}}\\
\quad Arrow of time & $\log_{10}(S_H)\approx 14.1$ & $>0$ & \textcolor{blue}{\textbf{Proxy (suite)}}\\
\quad GW background & $\Omega_{\rm gw}\approx 3\times10^{-17}$ & $\ll$ bounds & \textcolor{blue}{\textbf{Consistent (suite)}}\\
\quad Precision QED audit & TFPT-scale NP & $\ll$ anomaly scales & \textcolor{blue}{\textbf{Consistent (suite)}}\\
\midrule
\multicolumn{4}{l}{\textit{Requires Further Development}}\\
\quad CKM full matrix & Yukawa textures & RG-dressed baseline + $\chi^2$ diagnostic & \textcolor{blue}{\textbf{Implemented (suite; p-value gate satisfied)}}\\
\quad EW scale (constant factory) & $v$ (EW vev) & $320~\mathrm{GeV}$ vs $246~\mathrm{GeV}$ & \textcolor{orange}{\textbf{Pending (candidate; theoryv3 constant factory)}}\\
\quad Lepton ladder (constant factory) & $m_e, m_\mu, m_\tau$ & depends on $v$ candidate & \textcolor{orange}{\textbf{Pending (candidate chain)}}\\
\quad Hadronic scales & $m_p$ placeholder; $m_\pi$ pending & --- & \textcolor{orange}{\textbf{Open (not derived)}}\\
\quad Derived $H_0$ (diagnostic) & from $\Omega_b$ and $\eta_b$ proxy & $z_{\mathrm{Planck}}\simeq -2.99$ & \textcolor{blue}{\textbf{Diagnostic (open interface)}}\\
\quad $\Omega_b$ identity & $(4\pi-1)\beta_{\text{rad}}$ & $z=\num[round-mode=places,round-precision=3]{\TFPTOmegaBZ}$ & \textcolor{blue}{\textbf{Proxy (APS seam candidate; embedding open)}}\\
\end{longtable}
\endgroup

%=============================================================================
\section{What TFPT Does \emph{Not} Assume}
%=============================================================================

The following are \emph{not} required:
\begin{itemize}[leftmargin=2em]
    \item[\ding{55}] Slow-roll inflation (specific model)
    \item[\ding{55}] String landscape or anthropic selection
    \item[\ding{55}] Extra dimensions as free parameters
    \item[\ding{55}] Dark matter entirely from axions
    \item[\ding{55}] New particles beyond minimal axion sector
    \item[\ding{55}] Any fitted or adjusted parameters
\end{itemize}

%=============================================================================
\section{Conclusions}
%=============================================================================

TFPT demonstrates that the fine-structure constant $\alpha$ can be derived 
within an explicit set of eight minimal axioms rooted in topology, geometry, 
and quantum consistency. The key results are:

\begin{enumerate}
    \item The topological coupling $\cthree = 1/(8\pi)$ from 11D Chern--Simons 
        quantization and the geometric scale $\phiz = 1/(6\pi) + 3/(256\pi^4)$ 
        from M\"obius fiber reduction are the only fundamental inputs.
    
    \item The cubic fixed-point equation (CFE) has a unique positive root. With the derived two-defect term $\delta_2=\frac54(\delta_{\mathrm{top}})^2$ (holonomy $g=5$), this yields
        $\alpha^{-1}(0)=\num[round-mode=places,round-precision=12]{\TFPTalphaInvGOverFourNum}$ (one-defect truncation: $\num[round-mode=places,round-precision=10]{\TFPTalphaInvSelfNum}$).
    
    \item The same invariants determine the axion--photon coupling 
        $\gagg = -4\cthree$ and predict cosmic birefringence 
        $\beta = \num[round-mode=places,round-precision=4]{\TFPTbetaDegNum}^\circ$, consistent with observations.
    
    \item An E$_8$ cascade organizes all physical scales from Planck to 
        electroweak, with deviations typically below 5\%.
\end{enumerate}

The theory is falsifiable: any significant revision of $\alpha$, $\beta$, 
or the two-loop RG fingerprints would invalidate the framework.

\paragraph{Outlook (open interfaces).}
Two longer-range directions are sketched in internal notes and remain explicitly open:
an asymptotic-safety UV completion for the torsionful gravity sector and a black-hole
vortex/remnant picture with information-paradox implications. Both are framed as outlook
items rather than closed predictions at this stage \cite{fiveProblemsNote2025,blackHolesBh1Note2025}.

%=============================================================================
\appendix
\section{Numerical Constants (High Precision)}
%=============================================================================

\begin{table}[H]
\centering
\small
\begin{tabular}{@{}ll@{}}
\toprule
\textbf{Quantity} & \textbf{Value}\\
\midrule
$\cthree = 1/(8\pi)$ & \num[round-mode=places,round-precision=12]{\TFPTcThreeNum}\\
$\phitree = 1/(6\pi)$ & \num[round-mode=places,round-precision=12]{\TFPTvarphiTreeNum}\\
$\deltatop = 3/(256\pi^4)$ & \num[round-mode=places,round-precision=12]{\TFPTdeltaTopNum}\\
$\phiz = \phitree + \deltatop$ & \num[round-mode=places,round-precision=12]{\TFPTvarphiZeroNum}\\
$\phiz(\alpha)$ (self-consistent) & \num[round-mode=places,round-precision=12]{0.053170209119023745}\\
$\beta_{\text{rad}} = \phiz/(4\pi)$ & \num[round-mode=places,round-precision=12]{\TFPTbetaRadNum}\\
$\beta_{\text{deg}} = (180/\pi)\,\beta_{\text{rad}}$ & \num[round-mode=places,round-precision=6]{\TFPTbetaDegNum}$^\circ$\\
$g_{a\gamma\gamma}$ & \num[round-mode=places,round-precision=12]{\TFPTgAggNum}\\
$\gamma(0)=g/(g+1)$ (E$_8$) & $\tfrac56$\\
$\lambda$ (E$_8$ exponent) & \num[round-mode=places,round-precision=12]{0.587233190787837868}\\
$M/\Mpl$ ($R^2$ scale) & \num[round-mode=places,round-precision=12]{\TFPTMOverMplNum}\\
$A = 2\cthree^3$ & \num[round-mode=figures,round-precision=8]{1.2598255637968552e-4}\\
$\mathcal{K} = (\bone/2\pi)\ln(1/\phiz)$ & \num[round-mode=places,round-precision=10]{1.91468479457699546}\\
$\alpha^{-1}_{\text{baseline}}$ & \num[round-mode=places,round-precision=9]{137.036501464885822}\\
$\alpha^{-1}_{\text{TFPT}}$ (one-defect) & \num[round-mode=places,round-precision=12]{\TFPTalphaInvSelfNum}\\
$\alpha^{-1}_{\text{TFPT}}$ (two-defect, $g=5$) & \num[round-mode=places,round-precision=12]{\TFPTalphaInvGOverFourNum}\\
$\alpha^{-1}_{\text{TFPT}}$ (two-defect, suite diagnostic) & \num[round-mode=places,round-precision=12]{\TFPTalphaInvTwoDefNum}\\
$\alpha^{-1}_{\text{CODATA 2022}}$ & \num[round-mode=places,round-precision=9]{\TFPTalphaInvCODATANum}\\
Accuracy (ppm; one-defect) & \num[round-mode=places,round-precision=6]{-0.03704}\\
Accuracy (ppm; two-defect, $g=5$) & \num[round-mode=places,round-precision=6]{\TFPTalphaPpmGOverFourNum}\\
Accuracy (ppm; two-defect, suite diagnostic) & \num[round-mode=places,round-precision=6]{\TFPTalphaPpmTwoDefNum}\\
$\xi = \cthree/\phiz$ & \num[round-mode=places,round-precision=9]{0.748303083562547799}\\
\bottomrule
\end{tabular}
\caption{High-precision numerical values of TFPT constants.}
\label{tab:constants}
\end{table}

\begin{remark}
\textbf{Diagnostic only:} an exact CODATA match in the single-defect truncation would require
$\delta_2=\num[round-mode=figures,round-precision=6]{\TFPTdeltaTwoOptimalNum}$.
This value is listed only as a calibration diagnostic and is not treated as a prediction.
\end{remark}

%=============================================================================
\section{Diagnostic $g$ Variants (Internal Sanity Checks)}
\label{app:gdiagnostic}
%=============================================================================

The suite records legacy diagnostic templates in which $g$ is varied away from the holonomy
value $g=5$ (e.g.\ the ordered-occupancy $g=4$ template). These runs are used only to verify
that the defect-sector response is sensitive to discrete holonomy choices and to quantify the
direction of the shift in $\alpha^{-1}(0)$. They are explicitly \emph{not} used for parameter
selection or fitting. The audit trail is carried by \texttt{alpha\_precision\_audit} and the
independence of $g$ from any $\alpha$ reference inputs is enforced by the
\texttt{g5\_independence\_audit}.

%=============================================================================
\section{RG Unification Overview (Supplementary)}
%=============================================================================

\begin{figure}[H]
\centering
\begin{tikzpicture}
\begin{axis}[
    width=0.9\linewidth,
    height=0.48\linewidth,
    xlabel={$\log_{10}(\mu/\mathrm{GeV})$},
    ylabel={$\alpha_i^{-1}(\mu)$},
    title={Gauge-coupling unification overview},
    grid=both,
    legend pos=north east
]
\addplot table[x=log10mu, y expr={1/\thisrow{alpha1GUT}}] {\RgTable};
\addlegendentry{$\alpha_1^{-1}$}
\addplot table[x=log10mu, y expr={1/\thisrow{alpha2}}] {\RgTable};
\addlegendentry{$\alpha_2^{-1}$}
\addplot table[x=log10mu, y expr={1/\thisrow{alpha3}}] {\RgTable};
\addlegendentry{$\alpha_3^{-1}$}
\addplot[black, dashed] coordinates {(15.692222828950541,0) (15.692222828950541,200)};
\addlegendentry{$\mu_\ast\approx4.92\times10^{15}$ GeV}
\node[anchor=west] at (axis cs:15.9,110) {mismatch $\approx5.92$};
\end{axis}
\end{tikzpicture}
\caption{Gauge coupling convergence. Visualizes the unification gate scan used by the suite.}
\label{fig:rg-unification}
\end{figure}

%=============================================================================
\section{Coefficient Policy Sketch for the CFE}
%=============================================================================

This appendix records the coefficient bookkeeping used in the suite and provides a
heuristic motivation for the CFE coefficient policy. It is not presented as a full
background-field derivation. In the declared scheme, the one-loop effective potential
for the electromagnetic coupling modulus is taken to have the form:
\begin{equation}
    U(\alpha) = \frac{1}{4}(F_{\mu\nu}F^{\mu\nu})_{\text{eff}} 
        + \text{quantum corrections}
\end{equation}

The adopted structure, including the Chern--Simons contribution from 11D reduction, is:
\begin{equation}
    U(\alpha) = \frac{A}{4}\alpha^4 - \frac{2}{3}A\cthree^3\alpha^3 
        - 8A\bone\cthree^6 L\,\alpha + \mathcal{O}(\alpha^5)
    \label{eq:potential}
\end{equation}
where $L = \ln(1/\phiz)$ encodes the geometric scale.

The modulus-closure condition $\delta\Gamma_{\mathrm{eff}}/\delta\sigma=0$ reduces to
$\partial U/\partial\alpha = 0$ (monotonic $\alpha(\sigma)$), giving:
\begin{equation}
    \frac{\partial U}{\partial\alpha} = A\alpha^3 - 2A\cthree^3\alpha^2 
        - 8A\bone\cthree^6 L = 0
    \label{eq:stationarity}
\end{equation}

Dividing by $A$ (which is nonzero):
\begin{equation}
    \alpha^3 - 2\cthree^3\alpha^2 - 8\bone\cthree^6 L = 0
    \label{eq:CFE-appendix}
\end{equation}

This is the CFE under the declared policy. The combinatorial factors are motivated by
the following diagrammatic bookkeeping:

\textbf{(i) Factor 2 in $2\cthree^3\alpha^2$:}

The cubic term in the effective potential comes from the vacuum polarization 
(``bubble'') diagram with axion insertion:
\begin{center}
\begin{tikzpicture}[scale=0.8]
    \draw[thick, photon] (0,0) -- (1.5,0);
    \draw[thick] (1.5,0) arc (180:0:0.75);
    \draw[thick] (1.5,0) arc (-180:0:0.75);
    \draw[thick, photon] (3,0) -- (4.5,0);
    \draw[thick, dashed] (2.25,0.75) -- (2.25,1.5);
    \node at (2.25,1.8) {\small $a$};
    \node at (0,-0.3) {\small $\gamma$};
    \node at (4.5,-0.3) {\small $\gamma$};
\end{tikzpicture}
\end{center}
The $\varepsilon^{\mu\nu\rho\sigma}\varepsilon_{\alpha\beta\gamma\delta}$ contraction from two 
$F\tilde{F}$ vertices produces a factor 2 from the antisymmetric tensor algebra:
\begin{equation}
    \varepsilon^{\mu\nu\rho\sigma}\varepsilon_{\mu\nu\alpha\beta} = -2(\delta^\rho_\alpha\delta^\sigma_\beta - \delta^\rho_\beta\delta^\sigma_\alpha)
\end{equation}

\textbf{(ii) Factor 8 and $\cthree^6$ in $8\bone\cthree^6 L$:}

The log-dependent term has two components:

\textbf{Origin of $\cthree^6$:} The factor $\cthree^6$ arises from the product structure 
of the effective action. With $A = 2\cthree^3$ as the overall amplitude:
\begin{equation}
    8\bone\cthree^6 L = 8\bone \cdot \cthree^3 \cdot \cthree^3 \cdot L
    = 4 \cdot A \cdot \cthree^3 \cdot \bone L
\end{equation}
The factorization proceeds as follows:
\begin{itemize}[leftmargin=1.5em, topsep=2pt, itemsep=1pt]
    \item One factor of $\cthree^3$ comes from three $\cthree\, aF\tilde{F}$ vertices 
        in the anomaly chain (triangle diagram).
    \item The second factor of $\cthree^3$ comes from the amplitude normalization $A/2 = \cthree^3$.
    \item The coefficient 8 decomposes as $8 = 4 \times 2$, where 4 is from $|D_2|^2$ 
        (two $\mathbb{Z}_2$ reflections) and 2 is from the double cover sheets.
\end{itemize}
\textbf{Consistency check:} $8\cthree^6 = 8 \times (8\pi)^{-6} = 8/(8\pi)^6$, which matches 
the numerical coefficient in the CFE.

\textbf{Origin of factor 8:} The dihedral symmetry $D_4$ of the 4-point amplitude 
contributes $|D_4| = 8$. This arises from the 4 rotations and 4 reflections that 
leave the fermion loop topology invariant:
\begin{center}
\begin{tikzpicture}[scale=0.8]
    \draw[thick] (0,0) rectangle (2,2);
    \draw[thick, photon] (-0.8,0) -- (0,0);
    \draw[thick, photon] (-0.8,2) -- (0,2);
    \draw[thick, photon] (2,0) -- (2.8,0);
    \draw[thick, photon] (2,2) -- (2.8,2);
    \node at (1,-0.5) {\small fermion loop};
    \node at (1,1) {\small $D_4$};
\end{tikzpicture}
\end{center}

\textbf{The $\bone$ factor:} $\bone = \frac{3}{5}b_Y = 41/10$ is the GUT-normalized
hypercharge trace with explicit Weyl/scalar weights:
$b_Y=\frac{2}{3}\sum_{\text{Weyl}} d_3 d_2\,Y^2+\frac{1}{3}\sum_{\text{scalars}} d_3 d_2\,Y^2$.

\textbf{(iii) Factor 2/3 in potential:}

We choose the coefficient $2/3$ in the cubic term of \eqref{eq:potential} so that
differentiation produces the factor 2 in \eqref{eq:stationarity}:
\begin{equation}
    \frac{d}{d\alpha}\left(-\frac{2}{3}A\cthree^3\alpha^3\right) = -2A\cthree^3\alpha^2
\end{equation}
This choice is consistent with the integral representation used in the suite
bookkeeping: the cubic term is motivated by a triangle subgraph.

\textbf{Why no $\alpha^2$ term?}

The absence of an $\alpha^2$ term in the CFE is enforced by a declared renormalization
condition: the quadratic coefficient is absorbed into a counterterm fixed at the reference
point of the chosen scheme. The toy one-loop calculation shows that an $\alpha^2$ term is
generic, so the condition must be stated explicitly rather than implied by Ward identities.
The toy calculation is a structure check only (the toy $\beta$-prefactor mismatch is flagged),
not a Standard Model one-loop derivation. The resulting precision requirement is strong
($|\epsilon_2|\lesssim 1.58\times 10^{-10}$ for 1 ppm), and is treated as a formal constraint
of the kernel EFT policy (see \texttt{alpha2\_term\_robustness\_scan}).

\textbf{Vacuum stability:} The stationary point is required to be a local minimum,
$U''(\alpha_\star)>0$ with a finite margin; this is enforced by the
\texttt{vacuum\_stability\_audit} and recorded in the reproducibility appendix.

\begin{tcolorbox}[colback=purple!3, colframe=purple!50!black, 
    title={\textbf{Why Integer Coefficients? (No $\zeta(3)$, $\ln 2$, etc.)}}, fonttitle=\bfseries\small]
\textbf{Potential objection:} In typical QFT calculations (e.g., anomalous magnetic moment $g-2$), 
coefficients involve transcendental numbers like $\zeta(3)$, $\pi^2/12$, $\ln 2$. Why are the 
CFE coefficients (2, 8, 48) simple integers?

\textbf{Resolution:} The CFE is not a perturbative loop expansion but a \emph{topological constraint}.

\begin{enumerate}[leftmargin=1.5em, topsep=2pt, itemsep=1pt]
    \item \textbf{Origin of transcendentals:} In $g-2$, transcendentals arise from 
        \emph{integrating over continuous momenta} in Feynman integrals. The integrals 
        produce polylogarithms, which evaluate to $\zeta$-values.
    
    \item \textbf{Topological discreteness:} In TFPT, the relevant integrals are over 
        \emph{discrete topological data}: Euler characteristics, winding numbers, and 
        symmetry group orders. These are integers by definition.
    
    \item \textbf{Combinatorial origin:}
        \begin{itemize}[leftmargin=1em, topsep=1pt, itemsep=1pt]
            \item Factor 2: From $\varepsilon^{\mu\nu\rho\sigma}$ contraction (antisymmetry)
            \item Factor 8: From $|D_4| = 8$ (dihedral group of the box)
            \item Factor 48: \emph{Mnemonic} $2 \times 24$ (double cover $\times$ permutations); the derivation is the spin-lifted deficit with four cut patches (Appendix~\ref{app:48deriv})
        \end{itemize}
    
    \item \textbf{Where $\pi$ appears:} The factors of $\pi$ in $\cthree = 1/(8\pi)$ and 
        $\phitree = 1/(6\pi)$ come from the \emph{normalization} of circles ($2\pi$ per cycle), 
        not from loop integrals. They are geometric, not perturbative.
\end{enumerate}

\textbf{Diagnostic:} If transcendentals appeared in the CFE, it would indicate non-topological 
dynamics. Their absence is evidence that the fixed-point structure is genuinely topological.
\end{tcolorbox}

\begin{remark}
The consistency between \eqref{eq:potential} and \eqref{eq:CFE-appendix} can be 
verified directly: $\frac{d}{d\alpha}\bigl(-\frac{2}{3}A\cthree^3\alpha^3\bigr) 
= -2A\cthree^3\alpha^2$, matching the required coefficient.
\end{remark}

%=============================================================================
